%% Generated by verification/make_figures.py from regular-graph-harmonic-quotient-examples.tex.
%% Figure fig:harmonic-quotient_d3reg12n+twinC5 (replaces figures/harmonic-quotient_d3reg12n+twinC5-G2.tex).
%% Gamma_2 = d3reg12n: 12 vertices, 18 edges (E_2 in the example).
%% phi from the table at line 4735. Vertex x is filled with color hq<phi(x)>.
%% Layout: circle, Hamiltonian cycle 0,1,2,3,9,8,11,5,4,10,7,6 (9 crossings).
\begin{tikzpicture}[hqgraph]
\node[hqv, fill=hq2] (v0) at (0.000,1.770) {0};
\node[hqv, fill=hq1] (v1) at (-0.885,1.533) {1};
\node[hqv, fill=hq1] (v2) at (-1.533,0.885) {2};
\node[hqv, fill=hq2] (v3) at (-1.770,0.000) {3};
\node[hqv, fill=hq3] (v4) at (1.533,-0.885) {4};
\node[hqv, fill=hq3] (v5) at (0.885,-1.533) {5};
\node[hqv, fill=hq0] (v6) at (0.885,1.533) {6};
\node[hqv, fill=hq5] (v7) at (1.533,0.885) {7};
\node[hqv, fill=hq5] (v8) at (-0.885,-1.533) {8};
\node[hqv, fill=hq0] (v9) at (-1.533,-0.885) {9};
\node[hqv, fill=hq4] (v10) at (1.770,-0.000) {10};
\node[hqv, fill=hq4] (v11) at (-0.000,-1.770) {11};
\draw[hqe] (v0) -- (v1);
\draw[hqe] (v0) -- (v5);
\draw[hqe] (v0) -- (v6);
\draw[hqe] (v1) -- (v2);
\draw[hqe] (v1) -- (v7);
\draw[hqe] (v2) -- (v3);
\draw[hqe] (v2) -- (v8);
\draw[hqe] (v3) -- (v4);
\draw[hqe] (v3) -- (v9);
\draw[hqe] (v4) -- (v5);
\draw[hqe] (v4) -- (v10);
\draw[hqe] (v5) -- (v11);
\draw[hqe] (v6) -- (v7);
\draw[hqe] (v6) -- (v9);
\draw[hqe] (v7) -- (v10);
\draw[hqe] (v8) -- (v9);
\draw[hqe] (v8) -- (v11);
\draw[hqe] (v10) -- (v11);
\end{tikzpicture}%
